\chapter{The \textnormal{\texttt{show}} Command Cookbook}\label{app:show}

Eight of the twenty-nine blueprint topics begin \emph{Troubleshoot} or
\emph{Diagnose}, and topic~\tp{2.4} names \cmd{show} commands explicitly. This
appendix is organised by \textbf{the question you are asking}, because that is
what you have when a fault starts --- not by command name, which is what you
have afterwards.

Each entry gives the question, the command, and \textbf{the specific field to
read}. Running a command and skimming its output is not diagnosis; knowing which
line answers your question is.

\begin{conceptbox}[Four modifiers that make all of these faster]
\begin{tight}
  \item \cmd{| include <text>} --- only lines containing it.
  \item \cmd{| exclude <text>} --- everything except those lines.
  \item \cmd{| begin <text>} --- start output at the first match, and continue.
  \item \cmd{| section <text>} --- the whole configuration block matching it.
\end{tight}

\cmd{show running-config | section interface GigabitEthernet1/0/5} is the fastest
way to see one port's configuration, and \cmd{show interfaces status | exclude
connected} shows only the ports that are not working --- which on a 48-port
switch is the difference between reading four lines and reading forty-eight.
\end{conceptbox}

\section{Is the interface healthy? \textnormal{(\chref{physical})}}

\begin{longtable}{@{}L{4.4cm}L{5.6cm}L{5.3cm}@{}}
\toprule
\rowcolor{primary!15}
\textbf{Question} & \textbf{Command} & \textbf{Read this} \\
\midrule
\endhead
Which ports are up, and with what address? &
\cmd{show ip interface brief} &
\cmd{Status} and \cmd{Protocol}. \cmd{Method} says \cmd{DHCP} or \cmd{manual} \\
\rowcolor{lightbg}
Which ports are \emph{not} working? &
\cmd{show interfaces status | exclude connected} &
\cmd{notconnect}, \cmd{err-disabled}, \cmd{disabled} \\
Is the physical link clean? &
\cmd{show interfaces <int>} &
\cmd{input errors}, \cmd{CRC}, \cmd{late collision}, \cmd{runts}. CRC with no
collisions suggests cable or duplex \\
\rowcolor{lightbg}
Is duplex mismatched? &
\cmd{show interfaces <int> | include duplex} &
Half duplex on a switch-to-switch link is almost always wrong \\
Why is the port err-disabled? &
\cmd{show interfaces status err-disabled} &
The \cmd{Reason} column names the feature that did it \\
\rowcolor{lightbg}
Is the optic within spec? &
\cmd{show interfaces transceiver} &
Rx power against the stated range \\
\bottomrule
\end{longtable}

\section{Is layer 2 right? \textnormal{(\chref{vlans}--\chref{stp})}}

\begin{longtable}{@{}L{4.4cm}L{5.6cm}L{5.3cm}@{}}
\toprule
\rowcolor{primary!15}
\textbf{Question} & \textbf{Command} & \textbf{Read this} \\
\midrule
\endhead
Which VLAN is this port in? &
\cmd{show vlan brief} &
The port should be listed against the VLAN you expect \\
\rowcolor{lightbg}
Does the VLAN exist at all? &
\cmd{show vlan brief} &
A missing VLAN is a common cause of a dead access port \\
Is the trunk carrying what I think? &
\cmd{show interfaces trunk} &
\textbf{\cmd{Vlans allowed and active}} --- not the configured list. This is the
field that matters \\
\rowcolor{lightbg}
Do the native VLANs match? &
\cmd{show interfaces trunk} &
\cmd{Native vlan}, on both ends. A mismatch is logged by CDP \\
Is the EtherChannel up? &
\cmd{show etherchannel summary} &
\cmd{(P)} in the bundle, \cmd{(I)} independent, \cmd{(s)} suspended \\
\rowcolor{lightbg}
Who is the root bridge? &
\cmd{show spanning-tree vlan <id>} &
\cmd{This bridge is the root}, or the root's ID and the port cost to it \\
Which ports are blocking? &
\cmd{show spanning-tree vlan <id> | include BLK} &
Any unexpected blocking port \\
\rowcolor{lightbg}
Is there a loop? &
\cmd{show logging | include MACFLAP} &
A MAC flapping between two ports \\
Where is this MAC? &
\cmd{show mac address-table address <mac>} &
The port it was last seen on \\
\bottomrule
\end{longtable}

\section{What is attached? \textnormal{(\chref{discovery})}}

\begin{longtable}{@{}L{4.4cm}L{5.6cm}L{5.3cm}@{}}
\toprule
\rowcolor{primary!15}
\textbf{Question} & \textbf{Command} & \textbf{Read this} \\
\midrule
\endhead
What Cisco device is on this port? &
\cmd{show cdp neighbors} &
Local port, remote device, remote port \\
\rowcolor{lightbg}
What is its address and version? &
\cmd{show cdp neighbors detail} &
Management address, platform, IOS version \\
What non-Cisco device is attached? &
\cmd{show lldp neighbors} &
LLDP must be enabled first --- it is off by default \\
\bottomrule
\end{longtable}

\section{Where does traffic go? \textnormal{(\chref{routingtable}--\chref{fhrp})}}

\begin{longtable}{@{}L{4.4cm}L{5.6cm}L{5.3cm}@{}}
\toprule
\rowcolor{primary!15}
\textbf{Question} & \textbf{Command} & \textbf{Read this} \\
\midrule
\endhead
How will this packet be routed? &
\cmd{show ip route <destination>} &
The single entry that would be used --- longest prefix match already applied \\
\rowcolor{lightbg}
What does the table hold? &
\cmd{show ip route} &
Codes, then \cmd{[AD/metric]} on each entry \\
Is there a default route? &
\cmd{show ip route | include 0.0.0.0} &
\cmd{Gateway of last resort} \\
\rowcolor{lightbg}
Are OSPF neighbours up? &
\cmd{show ip ospf neighbor} &
\cmd{FULL} is healthy. \cmd{EXSTART} means \textbf{MTU mismatch} \\
Why will they not form? &
\cmd{show ip ospf interface <int>} &
Area, hello and dead timers, network type --- all must match \\
\rowcolor{lightbg}
Which OSPF interfaces are running? &
\cmd{show ip ospf interface brief} &
Cost, state, neighbour count \\
Is IPv6 routing working? &
\cmd{show ipv6 route ospf} &
Next hops are \cmd{fe80::} link-local, and that is correct \\
\rowcolor{lightbg}
Which router is the active gateway? &
\cmd{show standby brief} &
\cmd{Active}, \cmd{Standby}, priority, \cmd{P} for preempt \\
Same for VRRP? &
\cmd{show vrrp brief} &
\cmd{Master} or \cmd{Backup} \\
\bottomrule
\end{longtable}

\section{Are the services working? \textnormal{(\chref{dhcp}, \chref{nat}, \chref{acl})}}

\begin{longtable}{@{}L{4.4cm}L{5.6cm}L{5.3cm}@{}}
\toprule
\rowcolor{primary!15}
\textbf{Question} & \textbf{Command} & \textbf{Read this} \\
\midrule
\endhead
Did anyone get a lease? &
\cmd{show ip dhcp binding} &
Address, MAC, expiry \\
\rowcolor{lightbg}
Is the pool exhausted? &
\cmd{show ip dhcp pool} &
\cmd{Leased} against \cmd{Total addresses} \\
Is something using a pool address statically? &
\cmd{show ip dhcp conflict} &
Conflicts stay out of the pool until cleared \\
\rowcolor{lightbg}
Is NAT translating? &
\cmd{show ip nat translations} &
\cmd{Pro} column: \cmd{tcp} means PAT; \cmd{---} means static or dynamic \\
Is NAT even being attempted? &
\cmd{show ip nat statistics} &
\textbf{\cmd{Inside interfaces} and \cmd{Outside interfaces}}. Empty means the
markers are missing. \cmd{Hits: 0, Misses: 0} means no packet was considered \\
\rowcolor{lightbg}
Is the ACL dropping things? &
\cmd{show access-lists} &
The match counters. Zero on the last line means everything is being dropped;
zero in the middle means a shadowed rule \\
Is the ACL even applied? &
\cmd{show ip interface <int> | include access list} &
Inbound and outbound, named separately \\
\bottomrule
\end{longtable}

\section{What is the security state? \textnormal{(\chref{l2security})}}

\begin{longtable}{@{}L{4.4cm}L{5.6cm}L{5.3cm}@{}}
\toprule
\rowcolor{primary!15}
\textbf{Question} & \textbf{Command} & \textbf{Read this} \\
\midrule
\endhead
Is port security refusing something? &
\cmd{show port-security interface <int>} &
\cmd{Security Violation Count} on a \cmd{Secure-up} port usually means the
maximum is too low \\
\rowcolor{lightbg}
Which ports have violations? &
\cmd{show port-security} &
Violation counts across the switch \\
Is snooping building bindings? &
\cmd{show ip dhcp snooping binding} &
Empty for a device means it is statically addressed --- which breaks DAI \\
\rowcolor{lightbg}
Is DAI dropping ARP? &
\cmd{show ip arp inspection statistics vlan <id>} &
\cmd{DHCP Drops} means no binding; \cmd{ACL Drops} means an ARP ACL denied it \\
Is a port suppressing a storm? &
\cmd{show storm-control broadcast} &
\cmd{Blocking} means it is suppressing now \\
\bottomrule
\end{longtable}

\section{What happened? \textnormal{(\chref{management})}}

\begin{longtable}{@{}L{4.4cm}L{5.6cm}L{5.3cm}@{}}
\toprule
\rowcolor{primary!15}
\textbf{Question} & \textbf{Command} & \textbf{Read this} \\
\midrule
\endhead
\textbf{Why does the log show nothing?} &
\cmd{show logging | include level} &
\textbf{Run this before believing an empty log.} A threshold of \cmd{errors}
discards severity 4 to 7, which is most of what explains a fault \\
\rowcolor{lightbg}
What happened recently? &
\cmd{show logging} &
Newest at the bottom \\
Did somebody change something? &
\cmd{show logging | include CONFIG} &
\cmd{\%SYS-5-CONFIG\_I} --- somebody saved a change \\
\rowcolor{lightbg}
\textbf{Are there unsaved changes?} &
\cmd{show archive config differences} \newline
\cmd{~~nvram:startup-config} \newline
\cmd{~~system:running-config} &
\textbf{Any output at all means unsaved changes.} Make this the last step of
every change \\
How long has it been up? &
\cmd{show version} &
\cmd{uptime}, and the reason for the last reload \\
\rowcolor{lightbg}
Is the SNMP server reachable? &
\cmd{show aaa servers} \newline \cmd{show snmp} &
State and duration \\
\bottomrule
\end{longtable}

\section{The ten to know cold}

If you learn nothing else in this appendix, learn these. Between them they
answer the first question in almost every fault in this book.

\begin{longtable}{@{}L{7.4cm}L{7.9cm}@{}}
\toprule
\rowcolor{primary!15}
\textbf{Command} & \textbf{Answers} \\
\midrule
\endhead
\cmd{show ip interface brief} & Is it up, and what address does it have? \\
\rowcolor{lightbg}
\cmd{show interfaces status} & What is plugged in, and what is broken? \\
\cmd{show vlan brief} & Is this port in the VLAN I think? \\
\rowcolor{lightbg}
\cmd{show interfaces trunk} & Is the trunk carrying the VLAN? \\
\cmd{show ip route <destination>} & Where would this packet actually go? \\
\rowcolor{lightbg}
\cmd{show ip ospf neighbor} & Is the routing protocol healthy? \\
\cmd{show cdp neighbors} & What is on the other end of this cable? \\
\rowcolor{lightbg}
\cmd{show logging} & What happened, and when? \\
\cmd{show running-config | section <thing>} & What is this one thing configured
as? \\
\rowcolor{lightbg}
\cmd{show access-lists} & Is a filter responsible? \\
\bottomrule
\end{longtable}

\begin{alertbox}[The habit that makes all of this work]
Run these on a \emph{working} network first, and read the output until it is
familiar. You cannot recognise abnormal output if you have never looked at normal
output, and the exam will show you a \cmd{show} dump and expect you to spot the
one wrong field in it.

Every lab in this volume asks you to record normal output before breaking
anything. That instruction is what this appendix is for.
\end{alertbox}
